\section{مدول‌های تصویری و انژکتیو}

% Source: sources/lectures/17 Injective and Projective Modules.pdf | PDF page: 1 | item: projective module definition and lifting diagram
\begin{frame}[t]{تعریف مدول تصویری}
\label{lec17:projective-definition}
\begin{definition}
$R$-مدول $P$ را \textit{تصویری} یا \textit{پروژکتیو} می‌نامیم هرگاه برای هر
دو $R$-مدول $M,N$، هر $R$-همریختی $f:P\to N$ و هر دنبالهٔ دقیق
\[
M\xrightarrow{\varphi}N\longrightarrow0,
\]
یک $R$-همریختی $\psi:P\to M$ وجود داشته باشد به‌طوری‌که
$\varphi\psi=f$.
\end{definition}
\[
\begin{tikzcd}[ampersand replacement=\&,column sep=large,row sep=large]
\& P\arrow[d,"f"]\arrow[dl,dashed,"\psi"'] \& \\
M\arrow[r,"\varphi"']\&N\arrow[r]\&0
\end{tikzcd}
\]
\end{frame}

% Source: sources/lectures/17 Injective and Projective Modules.pdf | PDF page: 1 | item: Z is projective over Z
\begin{frame}[t]{تصویری‌بودن مدول اعداد صحیح}
\label{lec17:integers-projective}
\begin{example}
$\mathbb Z$ به‌عنوان $\mathbb Z$-مدول تصویری است.
\end{example}
\begin{proof}
فرض کنید $\varphi:M\to N$ پوشا و $f:\mathbb Z\to N$ همریختی باشد.
چون $f(1)\in N$، عضوی $m\in M$ هست که $\varphi(m)=f(1)$.
تعریف می‌کنیم
\[
\psi:\mathbb Z\longrightarrow M,\qquad \psi(n)=nm.
\]
برای هر $n\in\mathbb Z$ داریم $f(n)=nf(1)$ و
\[
\varphi\psi(n)=\varphi(nm)=n\varphi(m)=nf(1)=f(n).
\]
پس $\varphi\psi=f$.
\end{proof}
\end{frame}

% Source: sources/lectures/17 Injective and Projective Modules.pdf | PDF page: 2 | item: free modules are projective, basis lift and linear extension
\begin{frame}[t]{هر مدول آزاد تصویری است}
\label{lec17:free-is-projective}
\begin{theorem}
هر $R$-مدول آزاد، تصویری است.
\end{theorem}
\begin{proof}
فرض کنید $L$ آزاد و $\{x_i\}_{i\in I}$ پایهٔ آن باشد.
برای $\varphi:M\to N$ پوشا و $f:L\to N$، برای هر $i\in I$ عضوی
$m_i\in M$ انتخاب می‌کنیم که
\[
\varphi(m_i)=f(x_i).
\]
با قرار دادن $\psi(x_i)=m_i$ و گسترش خطی، تعریف می‌کنیم
\[
\psi\!\left(\sum_{i\in I}r_ix_i\right)=\sum_{i\in I}r_im_i
\qquad\text{(جمع‌ها متناهی‌اند).}
\]
آنگاه $\varphi\psi(x_i)=\varphi(m_i)=f(x_i)$؛ پس $\varphi\psi=f$.
\end{proof}
\end{frame}

% Source: sources/lectures/17 Injective and Projective Modules.pdf | PDF page: 2 | item: three equivalent characterizations of a projective module
\begin{frame}[t]{مشخصه‌های مدول تصویری}
\label{lec17:projective-characterizations}
\begin{theorem}
برای $R$-مدول $P$، شرط‌های زیر هم‌ارزند:
\begin{enumerate}
\item $P$ تصویری است.
\item هر دنبالهٔ دقیق کوتاهِ زیر شکافته می‌شود:
\[
0\longrightarrow M\xrightarrow{\varphi}N\xrightarrow{\psi}P
\longrightarrow0.
\]
\item یک $R$-مدول آزاد $F$ وجود دارد که برای یک $R$-مدول $A$،
\[
F\cong A\oplus P.
\]
\end{enumerate}
\end{theorem}
\end{frame}

% Source: sources/lectures/17 Injective and Projective Modules.pdf | PDF page: 3 | item: projectivity implies splitting via identity lift
\begin{frame}[t]{از تصویری‌بودن به شکافتن دنباله}
\label{lec17:projective-implies-split}
\begin{proof}[اثباتِ $\text{(۱)}\Rightarrow\text{(۲)}$]
در تعریف تصویری‌بودن، همریختی $\operatorname{id}_P:P\to P$ را در نظر بگیرید.
\[
\begin{tikzcd}[ampersand replacement=\&,column sep=large,row sep=large]
\&\&\&P\arrow[d,"\operatorname{id}_P"]\arrow[dl,dashed,"\theta"']\&\\
0\arrow[r]\&M\arrow[r]\&N\arrow[r,"\psi"']\&P\arrow[r]\&0
\end{tikzcd}
\]
همریختی $\theta:P\to N$ موجود است که
\[
\psi\theta=\operatorname{id}_P.
\]
پس دنباله شکافته می‌شود.
\end{proof}
\end{frame}

% Source: sources/lectures/17 Injective and Projective Modules.pdf | PDF page: 3 | item: construct free module indexed by elements of P and epimorphism to P
\begin{frame}[t]{ساخت مدول آزاد روی اعضای مدول}
\label{lec17:free-cover}
\begin{proof}[آغازِ اثباتِ $\text{(۲)}\Rightarrow\text{(۳)}$]
مجموعهٔ نمادهای $\mathcal X=\{x_p\}_{p\in P}$ را در نظر بگیرید و
$F$ را $R$-مدول آزاد با این پایه بگیرید:
\[
F=\left\{\sum_{p\in P}r_px_p\ \middle|\ r_p\in R,
\ \text{تنها متناهی ضریب ناصفر است}\right\}.
\]
تعریف $\varphi(x_p)=p$ با گسترش خطی، همریختیِ پوشای زیر را می‌دهد:
\[
\varphi:F\longrightarrow P,\qquad
\varphi\!\left(\sum_{p\in P}r_px_p\right)=\sum_{p\in P}r_pp.
\]
\end{proof}
\end{frame}

% Source: sources/lectures/17 Injective and Projective Modules.pdf | PDF page: 3 | item: kernel exact sequence and split free cover
\begin{frame}[t]{شکافتن پوشش آزاد}
\label{lec17:split-free-cover}
\begin{proof}[پایانِ اثباتِ $\text{(۲)}\Rightarrow\text{(۳)}$]
چون $\varphi:F\to P$ پوشاست، دنبالهٔ زیر دقیق است:
\[
0\longrightarrow\ker\varphi\longrightarrow F
\xrightarrow{\varphi}P\longrightarrow0.
\]
به فرضِ شرطِ (۲)، این دنباله شکافته می‌شود. بنابراین
\[
F\cong\ker\varphi\oplus P.
\]
با قرار دادن $A=\ker\varphi$، شرطِ (۳) به دست می‌آید.
\end{proof}
\end{frame}

% Source: sources/lectures/17 Injective and Projective Modules.pdf | PDF page: 4 | item: free direct summand implies projectivity, projection and inclusion diagram
\begin{frame}[t]{جمعوند مستقیمِ مدول آزاد}
\label{lec17:summand-implies-projective}
\begin{proof}[اثباتِ $\text{(۳)}\Rightarrow\text{(۱)}$]
$F\cong A\oplus P$ آزاد است. تصویر و شمول متناظر را به‌ترتیب
$\pi_P:F\to P$ و $i_P:P\to F$ بنامید.
برای $\varphi:M\to N$ پوشا و $f:P\to N$، چون $F$ تصویری است،
همریختی $\psi:F\to M$ موجود است که $\varphi\psi=f\pi_P$.
\[
\begin{tikzcd}[ampersand replacement=\&,column sep=large,row sep=large]
\&F\arrow[d,"\pi_P"]\arrow[ddl,bend right=15,"\psi"']\&\\
\&P\arrow[u,bend right=35,"i_P"']\arrow[d,"f"]\arrow[dl,dashed,"\psi_1"']\&\\
M\arrow[r,"\varphi"']\&N\arrow[r]\&0
\end{tikzcd}
\]
\end{proof}
\end{frame}

% Source: sources/lectures/17 Injective and Projective Modules.pdf | PDF page: 4 | item: restriction of lift and two corollaries
\begin{frame}[t]{تکمیل مشخصهٔ جمعوند مستقیم}
\label{lec17:summand-proof-conclusion}
با تعریف $\psi_1=\psi i_P$، برای هر $p\in P$ داریم
\[
\begin{aligned}
\varphi\psi_1(p)
&=\varphi\psi i_P(p)\\
&=\varphi\psi(0,p)\\
&=f\pi_P(0,p)=f(p).
\end{aligned}
\]
پس $P$ تصویری است.
\begin{corollary}
هر $R$-مدول تصویری، جمعوند مستقیمِ یک $R$-مدول آزاد است.
\end{corollary}
\begin{corollary}
هر جمعوند مستقیمِ یک $R$-مدول آزاد، تصویری است.
\end{corollary}
\end{frame}

% Source: sources/lectures/17 Injective and Projective Modules.pdf | PDF page: 5 | item: direct summand of a projective module is projective
\begin{frame}[t]{جمعوند مستقیمِ مدول تصویری}
\label{lec17:projective-summand-corollary}
\begin{corollary}
هر جمعوند مستقیمِ یک $R$-مدول تصویری، تصویری است.
\end{corollary}
\end{frame}

% Source: sources/lectures/17 Injective and Projective Modules.pdf | PDF page: 5 | item: Q is not projective over Z, basis coefficient proof
\begin{frame}[t]{اعداد گویا به‌عنوان مدول تصویری نیستند}
\label{lec17:rationals-not-projective}
\begin{example}
$\mathbb Q$ به‌عنوان $\mathbb Z$-مدول تصویری نیست.
\end{example}
\begin{proof}[آغاز برهان]
اگر تصویری باشد، برای مدول آزاد $F$ و مدول $A$ داریم
$A\oplus\mathbb Q\cong F$.
یک یکریختی $\theta:A\oplus\mathbb Q\to F$ و پایهٔ
$\{x_i\}_{i\in I}$ برای $F$ در نظر بگیرید. بنویسید
\[
\theta(0,1)=\sum_{i\in I}r_ix_i,\qquad r_i\in\mathbb Z.
\]
برای هر عدد صحیح مثبت $n$ نیز می‌توان نوشت
\[
\theta(0,1/n)=\sum_{i\in I}r'_ix_i.
\]
هر جمع تنها متناهی ضریب ناصفر دارد؛ ضرایبِ نوشته‌نشده را صفر می‌گیریم.
\end{proof}
\end{frame}

% Source: sources/lectures/17 Injective and Projective Modules.pdf | PDF page: 5 | item: all coefficients divisible by every n, contradiction, red marginal alternative
\begin{frame}[t]{برهانِ ناتصویری‌بودن اعداد گویا}
\label{lec17:rationals-coefficients}
\begin{proof}[ادامهٔ برهان]
برای هر $n\ge1$، از خطی‌بودن $\theta$ داریم
\[
\theta(0,1)=\theta\bigl(n(0,1/n)\bigr)
=n\theta(0,1/n)=\sum_i nr'_ix_i.
\]
یکتایی ضرایب نسبت به پایه می‌دهد
\[
r_i=nr'_i\quad\Longrightarrow\quad n\mid r_i
\qquad\text{برای هر }n\ge1.
\]
پس همهٔ $r_i$ها صفرند؛ در نتیجه $\theta(0,1)=0$،
که با یک‌به‌یک‌بودن $\theta$ تناقض دارد.
\end{proof}
\textbf{یادداشت حاشیهٔ جزوه.}
برای همریختی $\psi:\mathbb Q\to F$ با همین استدلال
$\psi(1)=0$ و سپس $\psi=0$؛ بنابراین شمولِ جمعوندِ $\mathbb Q$ ممکن نیست.
\end{frame}

% Source: sources/lectures/17 Injective and Projective Modules.pdf | PDF page: 5 | item: annihilating ideal allows quotient-ring module structure
\begin{frame}[t]{تغییر حلقه به خارج‌قسمت}
\label{lec17:quotient-ring-module}
\begin{remark}
اگر $M$ یک $R$-مدول و $I$ ایده‌آلی دوطرفه از $R$ باشد که
\[
I\subseteq\operatorname{Ann}_R(M),
\]
آنگاه $M$ یک $R/I$-مدول نیز هست.
\end{remark}
\end{frame}

% Source: sources/lectures/17 Injective and Projective Modules.pdf | PDF page: 5-6 | item: Z6 decomposition and projective nonfree modules; two conclusions
\begin{frame}[t]{مدول تصویریِ ناآزاد}
\label{lec17:z6-projective-nonfree}
% AMBIGUITY[L17-Q02]: primary, page 5 bottom; ring subscript and red quotient denominator unclear. Working reading 6 is provisional.
\begin{alertblock}{عدد حلقه با خوانش موقتِ \lr{L17-Q02}}
زیرنویس حلقه و مخرجِ کسر در دست‌نویس روشن نیست.
خوانش $6$ با تجزیهٔ نوشته‌شده سازگار است و اینجا موقتاً به کار می‌رود.
\end{alertblock}
می‌توان $\mathbb Z_2$ و $\mathbb Z_3$ را به‌عنوان
$\mathbb Z_6$-مدول در نظر گرفت. جزوه یادآوری می‌کند که
\[
\mathbb Z_6=\mathbb Z/6\mathbb Z,\qquad
\mathbb Z_6\cong\mathbb Z_2\oplus\mathbb Z_3.
\]
$\mathbb Z_6$ به‌عنوان $\mathbb Z_6$-مدول آزاد و در نتیجه تصویری است؛
پس هر دو جمعوند $\mathbb Z_2$ و $\mathbb Z_3$ تصویری‌اند، ولی آزاد نیستند.
\textbf{پرسش جزوه.} چرا این دو مدول آزاد نیستند؟
\end{frame}

% Source: sources/lectures/17 Injective and Projective Modules.pdf | PDF page: 6 | item: projective and free-summand nonfreeness conclusions following ambiguous example
\begin{frame}[t]{نتیجه‌های جزوه و شاهدِ موقت}
\label{lec17:nonfree-consequences}
% AMBIGUITY[L17-Q02]: these source conclusions follow the provisional cyclic example on page 5.
جزوه پس از مثال پیشین دو نتیجهٔ زیر را بیان می‌کند.
شاهدِ مثال برای این نتیجه‌ها تا روشن‌شدن عدد حلقه موقت است (\lr{L17-Q02}).
\begin{corollary}
هر $R$-مدول تصویری لزوماً آزاد نیست.
\end{corollary}
\begin{corollary}
هر جمعوندِ یک $R$-مدول آزاد لزوماً آزاد نیست.
\end{corollary}
\end{frame}

% Source: sources/lectures/17 Injective and Projective Modules.pdf | PDF page: 6 | item: direct sum of projectives iff every summand projective
\begin{frame}[t]{جمع مستقیمِ مدول‌های تصویری}
\label{lec17:projective-direct-sum}
\begin{theorem}
برای خانوادهٔ $\{P_i\}_{i\in I}$ از $R$-مدول‌ها، جمع مستقیم
$\bigoplus_{i\in I}P_i$ تصویری است اگر و تنها اگر هر یک از مدول‌های
$P_i$ تصویری باشد.
\end{theorem}
\begin{proof}[جهتِ نخست]
اگر $\bigoplus_{i\in I}P_i$ تصویری باشد، برای $F$ آزاد و مدول $A$ داریم
\[
\begin{aligned}
F&\cong\left(\bigoplus_{i\in I}P_i\right)\oplus A\\
&\cong P_j\oplus\left(\left(\bigoplus_{i\in I\setminus\{j\}}P_i\right)\oplus A\right).
\end{aligned}
\]
پس هر $P_j$ جمعوندی از یک مدول آزاد و بنابراین تصویری است.
\end{proof}
\end{frame}

% Source: sources/lectures/17 Injective and Projective Modules.pdf | PDF page: 6 | item: converse direct-sum proof and homomorphic-image remark
\begin{frame}[t]{جهتِ عکسِ قضیهٔ جمع مستقیم}
\label{lec17:projective-direct-sum-converse}
\begin{proof}
اگر هر $P_i$ تصویری باشد، مدول آزاد $F_i$ و مدول $A_i$ موجودند که
\[
F_i\cong P_i\oplus A_i.
\]
در نتیجه
\[
\bigoplus_{i\in I}F_i
\cong\left(\bigoplus_{i\in I}P_i\right)
\oplus\left(\bigoplus_{i\in I}A_i\right).
\]
مدول سمت چپ آزاد است؛ بنابراین $\bigoplus_{i\in I}P_i$ تصویری است.
\end{proof}
\begin{remark}
هر $R$-مدول تصویری، تصویر همریختِ یک $R$-مدول آزاد است.
\end{remark}
\end{frame}

% Source: sources/lectures/17 Injective and Projective Modules.pdf | PDF page: 7 | item: injective definition and extension diagram
\begin{frame}[t]{تعریف مدول انژکتیو}
\label{lec17:injective-definition}
\begin{definition}
$R$-مدول $J$ را \textit{انژکتیو} می‌نامیم هرگاه برای هر دنبالهٔ دقیق
$0\to M\xrightarrow{\varphi}N$ و هر $R$-همریختی $f:M\to J$،
یک $R$-همریختی $\psi:N\to J$ موجود باشد که
\[
\psi\varphi=f.
\]
\end{definition}
\[
\begin{tikzcd}[ampersand replacement=\&,column sep=large,row sep=large]
0\arrow[r]\&M\arrow[r,"\varphi"]\arrow[d,"f"']\&N\arrow[dl,dashed,"\psi"]\\
\&J\&
\end{tikzcd}
\]
\end{frame}

% Source: sources/lectures/17 Injective and Projective Modules.pdf | PDF page: 7 | item: Z is not injective, inclusion Z to Q, Hom(Q,Z)=0 argument
\begin{frame}[t]{اعداد صحیح انژکتیو نیستند}
\label{lec17:integers-not-injective}
\begin{example}
$\mathbb Z$ به‌عنوان $\mathbb Z$-مدول انژکتیو نیست.
\end{example}
\[
\begin{tikzcd}[ampersand replacement=\&,column sep=large,row sep=large]
0\arrow[r]\&\mathbb Z\arrow[r,hook]\arrow[d,"\operatorname{id}"']
\&\mathbb Q\arrow[dl,dashed,"\psi"]\\
\&\mathbb Z\&
\end{tikzcd}
\]
اگر $\psi\in\operatorname{Hom}_{\mathbb Z}(\mathbb Q,\mathbb Z)$ و
$\psi(1)=n$ باشد، برای هر عدد صحیح مثبت $m$ داریم
\[
n=\psi\!\left(m\frac1m\right)=m\psi\!\left(\frac1m\right).
\]
پس هر $m$ مقسوم‌علیهِ $n$ است و $n=0$؛ در نتیجه $\psi=0$.
بنابراین هیچ $\psi$ای نمی‌تواند همانیِ $\mathbb Z$ را گسترش دهد.
\end{frame}

% Source: sources/lectures/17 Injective and Projective Modules.pdf | PDF page: 8 | item: Baer criterion and forward implication
\begin{frame}[t]{محک بئر}
\label{lec17:baer-criterion}
\begin{theorem}
$R$-مدول $J$ انژکتیو است اگر و تنها اگر برای هر ایده‌آل چپ $I$ از $R$،
هر $R$-همریختی $f:I\to J$ را بتوان به $R$ گسترش داد.
\end{theorem}
\begin{proof}[جهتِ مستقیم]
اگر $J$ انژکتیو باشد، دنبالهٔ دقیق $0\to I\hookrightarrow R$ را در نظر بگیرید.
\[
\begin{tikzcd}[ampersand replacement=\&,column sep=large,row sep=large]
0\arrow[r]\&I\arrow[r,hook]\arrow[d,"f"']\&R\arrow[dl,dashed,"g"]\\
\&J\&
\end{tikzcd}
\]
آنگاه همریختی $g:R\to J$ موجود است که $g|_I=f$.
\end{proof}
\end{frame}

% Source: sources/lectures/17 Injective and Projective Modules.pdf | PDF page: 8-9 | item: reverse Baer proof, set of partial extensions and nonemptiness
\begin{frame}[t]{گسترش‌های جزئی در اثبات محک بئر}
\label{lec17:baer-partial-extensions}
فرض کنید شرط گسترش از ایده‌آل‌های چپ برقرار است. برای دنبالهٔ دقیق
\[
0\longrightarrow M\xrightarrow{\varphi}N
\]
و همریختی $f:M\to J$، مجموعهٔ زیر را در نظر بگیرید:
\[
\Sigma=\left\{h:M'\to J\ \middle|\
\operatorname{Im}\varphi\le M'\le N,\ h\in\operatorname{Hom}_R(M',J),\ h\varphi=f
\right\}.
\]
چون $\varphi:M\to\operatorname{Im}\varphi$ یکریختی است، همریختی
\[
f\varphi^{-1}:\operatorname{Im}\varphi\longrightarrow J
\]
در $\Sigma$ قرار دارد. پس $\Sigma\ne\varnothing$.
\end{frame}

% Source: sources/lectures/17 Injective and Projective Modules.pdf | PDF page: 9 | item: extension partial order and upper bound for a chain
\begin{frame}[t]{ترتیب گسترش و کران بالای زنجیر}
\label{lec17:baer-chain-upper-bound}
برای $h_1:M'_1\to J$ و $h_2:M'_2\to J$ در $\Sigma$ تعریف می‌کنیم
\[
h_1\le h_2\quad\Longleftrightarrow\quad
M'_1\le M'_2\ \text{و}\ h_2|_{M'_1}=h_1.
\]
یک زنجیر $\{h_i\}_{i\in I}$ در $\Sigma$ را در نظر بگیرید. تعریف می‌کنیم
\[
h:\bigcup_{i\in I}M'_i\longrightarrow J,\qquad
h(x)=h_i(x)\quad(x\in M'_i).
\]
چون اعضای زنجیر روی دامنه‌های مشترک توافق دارند، $h$ خوش‌تعریف است.
این $R$-همریختی در $\Sigma$ قرار دارد و کران بالای زنجیر است.
\end{frame}

% Source: sources/lectures/17 Injective and Projective Modules.pdf | PDF page: 9-10 | item: Zorn maximal extension and left ideal I
\begin{frame}[t]{عنصر ماکسیمال و ایده‌آل وابسته به آن}
\label{lec17:baer-maximal-extension}
به لم زورن، $\Sigma$ عنصری ماکسیمال مانند
\[
\widetilde h:\widetilde M\longrightarrow J
\]
دارد. اگر $\widetilde M=N$، گسترش مطلوب به دست آمده است.
در غیر این صورت، $n\in N\setminus\widetilde M$ انتخاب و تعریف می‌کنیم
\[
I=\{r\in R\mid rn\in\widetilde M\}.
\]
$0\in I$، پس $I\ne\varnothing$؛ همچنین $I$ یک ایده‌آل چپ از $R$ است.
\end{frame}

% Source: sources/lectures/17 Injective and Projective Modules.pdf | PDF page: 10 | item: define theta on I and extend it to R
\begin{frame}[t]{گسترش از ایده‌آل چپ}
\label{lec17:baer-ideal-extension}
تعریف می‌کنیم
\[
\theta:I\longrightarrow J,\qquad \theta(r)=\widetilde h(rn).
\]
چون برای $r\in I$ داریم $rn\in\widetilde M$، این نگاشت خوش‌تعریف و
یک $R$-همریختی است.
به فرضِ محک بئر، $\theta$ را می‌توان به همریختی‌ای مانند
\[
\widetilde\theta:R\longrightarrow J,\qquad
\widetilde\theta|_I=\theta
\]
گسترش داد.
\end{frame}

% Source: sources/lectures/17 Injective and Projective Modules.pdf | PDF page: 10-11 | item: define larger-domain extension, provisional scalar argument in formula
\begin{frame}[t]{ساخت گسترش روی دامنهٔ بزرگ‌تر}
\label{lec17:baer-larger-extension}
% AMBIGUITY[L17-Q01]: primary, page 10 lower formula has theta-tilde(rn); page 11 uses theta-tilde(r). Provisional scalar argument below.
\begin{alertblock}{خوانش موقتِ فرمول \lr{L17-Q01}}
صفحهٔ ۱۰ در آرگومانِ $\widetilde\theta$ عبارت $rn$ دارد؛
صفحهٔ ۱۱ از $r$ استفاده می‌کند. در ادامه، خوانش دوم موقتاً به کار می‌رود.
\end{alertblock}
با این خوانش، نگاشت زیر را در نظر می‌گیریم:
\[
\overline h:\widetilde M+\langle n\rangle\longrightarrow J,
\qquad
\overline h(m'+rn)=\widetilde h(m')+\widetilde\theta(r).
\]
برای $m'\in\widetilde M$ داریم
\[
\overline h(m')=\overline h(m'+0n)=\widetilde h(m'),
\qquad \overline h|_{\widetilde M}=\widetilde h.
\]
\end{frame}

% Source: sources/lectures/17 Injective and Projective Modules.pdf | PDF page: 11 | item: well-definedness computation of larger-domain extension
\begin{frame}[t]{خوش‌تعریفی گسترش؛ با خوانش موقت}
\label{lec17:baer-well-defined}
% AMBIGUITY[L17-Q01]: dependent proof, provisional theta-tilde(r) formula from page 11.
\textbf{ادامه با خوانش موقتِ \lr{L17-Q01}.}
اگر $m'_1+r_1n=m'_2+r_2n$، آنگاه
\[
(r_1-r_2)n=m'_2-m'_1\in\widetilde M,
\qquad r_1-r_2\in I.
\]
بنابراین
\[
\begin{aligned}
\widetilde h(m'_2)-\widetilde h(m'_1)
&=\widetilde h\bigl((r_1-r_2)n\bigr)\\
&=\theta(r_1-r_2)=\widetilde\theta(r_1-r_2)\\
&=\widetilde\theta(r_1)-\widetilde\theta(r_2).
\end{aligned}
\]
پس
\[
\widetilde h(m'_1)+\widetilde\theta(r_1)
=\widetilde h(m'_2)+\widetilde\theta(r_2),
\]
و $\overline h$ خوش‌تعریف است.
\end{frame}

% Source: sources/lectures/17 Injective and Projective Modules.pdf | PDF page: 11 | item: linearity, membership in Sigma, maximality contradiction and conclusion
\begin{frame}[t]{پایان برهان محک بئر؛ با خوانش موقت}
\label{lec17:baer-maximality-contradiction}
% AMBIGUITY[L17-Q01]: conclusion remains provisional pending formula review.
\textbf{ادامه با خوانش موقتِ \lr{L17-Q01}.}
چون $\widetilde h$ و $\widetilde\theta$ هر دو $R$-همریختی‌اند،
$\overline h$ نیز $R$-همریختی است. همچنین
\[
\widetilde M\lneq\widetilde M+\langle n\rangle.
\]
از $\operatorname{Im}\varphi\le\widetilde M$، برای هر $m\in M$ داریم
\[
\overline h\varphi(m)
=\overline h(\varphi(m))
=\widetilde h(\varphi(m))=f(m).
\]
بنابراین $\overline h\in\Sigma$ و گسترشی با دامنهٔ بزرگ‌تر از
$\widetilde h$ است؛ این با ماکسیمال‌بودن $\widetilde h$ تناقض دارد.
پس $\widetilde M=N$ و گسترش مطلوب حاصل می‌شود.
\end{frame}
